\chapter{Directed paths: the Cheng--Keevash results}
\label{ch:paths}

Thomass\'e's path conjecture, in the oriented form studied by Cheng
and Keevash (``Conjecture 1''), asserts that every oriented digraph
with minimum out-degree $\delta^+ \ge k$ contains a directed simple
path of length $2k$. The library proves the cases $k = 2$ and $k = 3$
and the general Lemma 7 of Cheng--Keevash, uniform in $\delta$.

\begin{trustbox}
\noindent\textbf{Trust.} All results axiom-free
(\rocqinline@Print Assumptions@ closed). Hypotheses are stated on
arbitrary finite oriented digraphs (\cref{def:oriented}) --- no
hidden strong-connectivity or regularity assumptions; where the
\emph{proof} reduces to a strong subdigraph, that reduction is part of
the formalized proof, not of the statement.
\end{trustbox}

\resulthead{The $\delta = 3$ path theorem}{ck_conj1_delta3}

\begin{lstlisting}[language=Rocq]
Theorem ck_conj1_delta3 (D : orientedDigraph) :
  0 < #|D| -> (forall v : D, 3 <= outdeg v) -> 6 <= ell D.
\end{lstlisting}
\rocqsource{theories/applications/ck3/ck3_main.v}{147}

\begin{lstlisting}[language=Rocq]
Corollary ck_conj1_delta3_path (D : orientedDigraph) :
  0 < #|D| -> (forall v : D, 3 <= outdeg v) ->
  exists x : D, exists s : seq D, dipath x s /\ size s = 6.
\end{lstlisting}
\rocqsource{theories/applications/ck3/ck3_main.v}{158}

\label{res:ck_conj1_delta3_path}

\begin{decoded}
Let $D$ be any nonempty finite oriented digraph in which every vertex
has out-degree at least $3$. Then $\ell(D) \ge 6$
(\cref{def:ell}): $D$ contains a directed simple path with $6$ arcs
(7 distinct vertices). The second form unfolds $\ell$ into the
explicit existence of such a path (\cref{def:dipath}).
This is the $k = 3$ case of Cheng--Keevash Conjecture~1.
\end{decoded}

\input{../web/panels/ck_conj1_delta3}

\resulthead{The $\delta = 2$ case}{ck_conj1_delta2}

\begin{lstlisting}[language=Rocq]
Corollary ck_conj1_delta2 (D : orientedDigraph) :
  0 < #|D| -> (forall v : D, 2 <= outdeg v) -> 4 <= ell D.
\end{lstlisting}
\rocqsource{theories/applications/ck3/lemma7.v}{698}


\begin{decoded}
Same statement one level down: minimum out-degree $\ge 2$ forces a
directed simple path of length $4$.
\end{decoded}

\begin{closurepanel}
\item \hyperref[def:arc]{Digraphs and arcs (u --> v)} (\S\ref*{def:arc})
\item \hyperref[def:oriented]{Oriented digraphs (irreflexive + asymmetric)} (\S\ref*{def:oriented})
\item \hyperref[def:outdeg]{Out-degree} (\S\ref*{def:outdeg})
\item \hyperref[def:dipath]{Directed simple paths} (\S\ref*{def:dipath})
\item \hyperref[def:ell]{The longest-dipath length ell(D)} (\S\ref*{def:ell})
\item plus the standard finite-mathematics vocabulary of the Glossary (\cref{ch:glossary})
\end{closurepanel}


\resulthead{Cheng--Keevash Lemma 7, uniform in $\delta$}{lemma7}

\begin{lstlisting}[language=Rocq]
Theorem lemma7 (D : orientedDigraph) (delta : nat) :
  0 < #|D| -> (forall v : D, delta <= outdeg v) ->
  2 * delta <= ell D \/
  exists S : {set D},
    [/\ S != set0, #|S| <= delta
      & forall v, v \in S -> 2 * delta - ell D <= outdeg_in S v].
\end{lstlisting}
\rocqsource{theories/applications/ck3/lemma7.v}{646}


\begin{decoded}
In a strong (\cref{def:strong}) oriented digraph that is
$\delta$-out-regular-from-below ($d^+(v) \ge \delta$ everywhere,
$\delta \ge 1$) and where the longest path is short
($\ell(D) \le 2\delta$), the vertex set of any longest path supports
the structural ``kernel'' conclusion quantified in the statement ---
the engine behind the path theorems. (This page quotes the statement
in full; the surrounding section fixes the oriented digraph
\rocqinline@D@, $\delta$, and the strongness/degree hypotheses
visible in the quoted text.)
\end{decoded}

\input{../web/panels/lemma7}

\resulthead{Theorem 4, oriented version}{ck_theorem4_oriented}

\begin{lstlisting}[language=Rocq]
Theorem ck_theorem4_oriented (D : orientedDigraph) (delta : nat) :
  0 < #|D| -> (forall v : D, delta <= outdeg v) ->
  2 * delta - (delta - 1)./2 <= ell D.
\end{lstlisting}
\rocqsource{theories/applications/ck3/lemma7.v}{684}


\begin{decoded}
Every nonempty oriented digraph with minimum out-degree $\ge \delta$
satisfies $\ell(D) \ge 2\delta - \lfloor(\delta-1)/2\rfloor$
(\rocqinline@./2@ is truncated halving). For even $\delta$ this is
sharper than the $\lceil 3\delta/2 \rceil$ stated by Cheng--Keevash.
\end{decoded}

\input{../web/panels/ck_theorem4_oriented}

\resulthead{No short-path strong counterexample at $\delta = 3$}{no_short_strong3}

\begin{lstlisting}[language=Rocq]
Lemma no_short_strong3 (H : orientedDigraph) :
  (forall v : H, outdeg v = 3) -> 0 < #|H| -> strongb H ->
  ell H <= 5 -> False.
\end{lstlisting}
\rocqsource{theories/applications/ck3/ck3_main.v}{28}


\begin{decoded}
There is no strong oriented digraph with $\delta^+ \ge 3$ and
$\ell \le 5$: the reduction target of the main theorem, stated
separately.
\end{decoded}

\input{../web/panels/no_short_strong3}
